\bookchapter{第三十二章：手機安全急救手冊}

\chapterintro{呢一章係一份「應急預案」。希望你永遠用唔上，但如果真係遇到手機丟失、被盜、被詐騙嘅情況，呢份手冊可以幫你將損失降到最低；地震、暴雨等災害嚟臨時，手機同樣係接收預警、定位求救、保持聯絡嘅工具。建議你而家就通讀一次，然後將關鍵步驟截圖保存，或者打印出嚟放喺屋企——因為出事嘅時候你嘅手機可能已經唔喺手上面，災害時就可能冇電或者冇訊號。}

\booksection{場景零：地震 / 極端天氣——手機點樣幫你保命}

防嘅係人身安全，唔係帳號資金。預警設定詳見 第十三章「天氣預警、自然災害預警同地震預警」；呢度只講警報響起之後應該點做。

\textbf{地震警報響咗：先躲後撤，唔好只顧住拍片}\par

\begin{enumerate}
  \item 就近避險：趴低、掩護、揸實（枱下、承重牆角、堅固傢俬旁），護住頭頸；遠離玻璃、外牆、吊燈。
  \item 唔好盲目衝樓：主震往往只有十幾秒到幾十秒，樓梯間可能更危險；電梯絕對唔可以用。
  \item 震停之後再有序撤離：行樓梯、留意餘震同墮物；去開闊地，遠離電線桿、廣告牌、危樓。
  \item 手機用喺刀刃上：畀緊急聯絡人發簡短訊息（「我喺 XX，安全/受傷，需要幫助」）比打長電話更慳電；可以發短訊/微信時優先文字。
  \item 冇訊號時：見 第七章 衛星 SOS / 衛星訊息（機型支援時）；或者去高處、開闊地再試；部分場景下無網藍牙對講亦有幫助。
  \item 關閉非必要功能：應急時開飛行模式再定時開、關藍牙掃描、降亮度，可以多撐幾小時；將最後一格電留畀報平安同救援訊息。
\end{enumerate}

\begin{tcbwarn}{暴雨 / 颱風 / 山洪：預警階段就要行動}
\begin{itemize}
  \item 睇到紅色預警：盡量減少外出，遠離地下空間、下穿橋、河道、陡坡。
  \item 喺戶外被困：向堅固建築物高處行，唔好強行涉水（積水可能有井蓋缺失同漏電）；打 119/110 時講清楚地標同水位。
  \item 出行前：用天氣 App 同當地預警公眾號確認風險；長途自駕避開山區臨水臨崖路段。
  \item 家人互報：約定「若失聯，以最後微信/短訊定位為準」，避免各自瘋狂撥號耗光電量。
\end{itemize}

\end{tcbwarn}

\begin{tcbpractice}{平時就做好嘅 4 件事}
\begin{enumerate}
  \item 開啟系統地震預警 / 正規預警 App，並允許佢突破靜音。
  \item 喺通訊錄置頂 2-3 位「緊急聯絡人」，鎖屏介面可以顯示（iOS/Android 醫療急救卡或者鎖屏資訊）。
  \item 屋企同袋入面備充電寶；長途出行前充滿電並帶線。
  \item 記住至少一個「唔用通訊錄都打得通」嘅號碼（家人手機號），寫喺紙上放銀包。
\end{enumerate}

\end{tcbpractice}

\booksection{場景一：手機唔見咗/俾人偷咗——黃金 30 分鐘}

手機丟失後嘅前 30 分鐘係「黃金時間」，處理得當可以保住你嘅錢財同私隱。按以下順序操作，一步都唔好跳：

\textbf{丟失後 30 分鐘標準操作流程}\par

\begin{enumerate}
  \item 借人哋手機撥打自己號碼（第 1 分鐘）：可能只係跌咗喺附近，好心人接咗就可以拎返嚟。
  \item 遙距鎖定手機（第 2-5 分鐘）：用人哋嘅手機或者電腦登入查找服務： iPhone：iCloud.com > 查找 > 揀你嘅裝置 > 標記做遺失
  \item Android 各品牌：華為 cloud.huawei.com / 小米 i.mi.com / OPPO cloud.oppo.com 等
  \item 掛失 SIM 卡（第 5-10 分鐘）：打電信商客服電話（移動 10086/聯通 10010/電信 10000），報手機丟失要求暫停 SIM 卡服務。呢一步最關鍵——騙子拎到你嘅 SIM 卡就可以接收驗證碼，重設你所有密碼。
  \item 凍結支付帳戶（第 10-20 分鐘）： 支付寶：95188 掛失，或者喺人哋手機上登入支付寶 > 我嘅 > 設定 > 安全中心 > 緊急凍結
  \item 微信：110 報警之後聯絡微信客服，或者喺微信安全中心 > 凍結帳號
  \item 銀行 App：撥打銀行卡背面嘅客服電話，要求凍結網銀
  \item 通知親友（第 20-25 分鐘）：喺朋友圈或者叫家人幫手發訊息，通知親友手機唔見咗，防止騙子用你嘅身份呃人。
  \item 報警（第 25-30 分鐘）：撥打 110 或者去最近嘅派出所報案，攞到報案回執。後續保險理賠、銀行糾紛都需要呢個回執。
\end{enumerate}

\booksection{場景二：俾人呃咗——72 小時黃金處理期}

發現自己俾人呃咗之後，最關鍵係冷靜。慌亂中做錯誤決定會導致二次損失。按以下步驟操作：

\begin{enumerate}
  \item 即刻停止一切轉帳操作：唔理騙子講乜嘢，唔好再轉一分錢。「再轉一筆保證金就可以退回之前嘅錢」係經典二次詐騙話術。
  \item 保留所有證據：唔好刪聊天記錄！唔好刪通話記錄！截圖保存聊天記錄、轉帳記錄、對方帳號資訊。呢啲都係報案同追回資金嘅關鍵證據。
  \item 即刻撥打 96110：呢個係全國反詐專線。接線員會指導你做緊急止付——如果轉帳時間唔長（通常 24 小時內），銀行有可能凍結對方帳戶攔截資金。
  \item 撥打 110 報警：帶齊手機、身份證、轉帳憑證去最近嘅派出所做筆錄。
  \item 聯絡銀行/支付平臺：告知俾人呃咗嘅情況，要求凍結涉案帳戶。如果係信用卡俾人盜刷，要求銀行啟動拒付流程。
  \item 修改所有密碼：喺另一部安全裝置上修改你嘅 Apple ID/品牌帳號、微信、支付寶、銀行 App 嘅密碼，開啟雙重認證。
\end{enumerate}

\booksection{場景三：手機中咗病毒/木馬——點算好？}

如果你發現手機出現以下症狀——莫名其妙彈出廣告、自動安裝奇怪 App、話費暴增、流量異常消耗——好可能係中咗惡意軟件：

\begin{itemize}
  \item 第一步：入安全模式——Android 手機長按電源鍵關機，再長按電源鍵開機，開機震動後長按音量減鍵唔放，直到入安全模式（左下角會顯示「安全模式」字樣）。安全模式下所有第三方 App 被停用，方便排查。
  \item 第二步：卸載可疑 App——喺安全模式下入設定 > 應用管理，按時間排序，搵到最近安裝嘅唔認識嘅 App 卸載佢。
  \item 第三步：安裝安全 App 掃描——下載「騰訊手機管家」或者「360手機衛士」做一次全盤掃描。
  \item 第四步：復原出廠設定（終極方案）——如果以上方法無效，備份重要數據後復原出廠設定。留意：復原出廠會刪除所有數據，一定要先備份相片、通訊錄等重要檔案。
  \item iPhone 基本唔會中毒——iOS 嘅沙盒機制令惡意軟件好難生存。如果你嘅 iPhone 出現異常，大概率係某個 App 嘅 bug 或者設定問題，唔係病毒。復原出廠設定可以解決 99\% 嘅問題。
\end{itemize}

\booksection{場景四：帳號被盜——快速奪返控制權}

\begin{enumerate}
  \item 即刻修改密碼：如果你仲可以登入帳號，第一時間改密碼。新密碼要同舊密碼完全不同。
  \item 踢走其他裝置：喺帳號安全設定入面搵到「登入裝置管理」，將唔認識嘅裝置全部踢下線。
  \item 檢查登入記錄：查看最近 7 日嘅登入記錄，如果有異地登入，就表示帳號真係被盜咗。
  \item 開啟雙重認證：如果你仲未開雙重認證，而家即刻開。呢個係防止帳號再次被盜最有效嘅方法。
  \item 檢查綁定資訊：檢查帳號綁定嘅手機號、電郵有冇被改。如果被改咗，用帳號申訴功能搵返。
  \item 檢查支付設定：如果被盜嘅係微信/支付寶，檢查免密支付、自動扣款有冇被篡改，銀行卡有冇被加到陌生帳戶。
\end{enumerate}

\booksection{場景五：收到勒索訊息——唔使驚，唔好畀錢}

常見嘅勒索套路有兩種：(1)「我有你嘅不雅影片/相片，唔畀錢就發畀你通訊錄入面所有人」；(2)「你嘅手機俾黑客控制咗，唔畀錢就鎖死」。遇到呢種情況：

\begin{itemize}
  \item 99.9\% 係詐騙：佢哋手入面根本冇你嘅任何私隱數據。羣發同樣嘅勒索訊息畀幾萬人，總有人驚到畀錢。
  \item 唔好回覆，唔好畀錢：回覆咗等於話畀對方知「呢個號碼有人用」，你會被標記做「可勒索目標」，後續騷擾更多。
  \item 截圖保存證據：然後拉黑對方。
  \item 修改密碼：如果擔心，改一下你主要帳號嘅密碼，求個安心。
  \item 如果真係有私隱洩露風險：報警。警方有處理呢類案件嘅經驗。
\end{itemize}

\booksection{場景六：AI 換臉/AI 擬聲詐騙——眼見不再為實}

由 2024 年起，AI 換臉同 AI 擬聲技術已經可以令騙子喺視像通話入面「變成」你嘅家人、上司、朋友嘅面同聲音。「媽，我出事啦，快啲轉錢過嚟」——影片入面嘅人長得、聽起嚟都同你個仔一模一樣，但其實係假嘅。呢個係目前最難防、亦最危險嘅詐騙手法之一。

\textbf{第一步：識別破綻——AI 偽造嘅四個常見漏洞}\par

\begin{itemize}
  \item 眨眼頻率異常：真人每分鐘眨眼 15-20 次，AI 換臉往往眨眼過少或者節奏僵硬。
  \item 嘴型同音訊唔同步：仔細睇口型同聽到嘅聲音之間，會有零點幾秒嘅錯位。
  \item 面部邊緣模糊：塊面嘅輪廓、耳仔附近、頭髮邊緣容易出現「毛邊」或者閃爍。
  \item 光影方向唔一致：塊面上面嘅高光、陰影同背景嘅光照方向對唔上，或者側面時五官「跟住轉」得唔啱。
\end{itemize}

\begin{tcbinfo}{第二步：三招驗證對方真假}
\begin{itemize}
  \item 要求做特定動作：叫對方轉頭、掩面、用手遮住個嘴再做開合動作。AI 換臉喺遮擋同大角度轉頭時最容易「穿煲」——塊面會變形或者跟唔上動作。
  \item 問只有對方先知嘅問題：編一個假問題，例如「上次我哋隻狗送去邊間醫院睇嘅？」——如果係騙子，佢答唔到或者支支吾吾。
  \item 回撥已知號碼：掛斷視像，用你通訊錄入面存嘅號碼直接回撥過去。騙子控制唔到真實嘅手機號，一撥就穿煲。
\end{itemize}

\end{tcbinfo}

\begin{tcbwarn}{第三步：確認俾人呃咗之後嘅報警流程}
\begin{enumerate}
  \item 即刻停止轉帳：再「急」都先停一下，凡係催你「快啲轉帳、唔好掛電話」嘅，幾乎都係騙子。
  \item 保留全部證據：錄屏、截圖保存視像通話畫面、聊天記錄、對方帳號、轉帳憑證——呢啲係技術鑑定嘅關鍵材料。
  \item 撥打 96110 反詐專線：告知係「AI 換臉/擬聲」類詐騙，接線員會啟動緊急止付。
  \item 撥打 110 報警：帶齊證據去派出所做筆錄，特別說明使用咗 AI 偽造技術，方便警方立案定性。
  \item 通知親友：騙子可能用你嘅偽造影片去呃人，提前喺朋友圈或者家庭羣入面發個提醒。
\end{enumerate}

\end{tcbwarn}

\booksection{場景七：SIM swap 攻擊（SIM 劫持）——你嘅號碼俾人「偷」走咗}

呢個係一種好多人未聽過、但危害極大嘅攻擊。SIM swap（SIM 卡劫持）指嘅係攻擊者透過偽造你嘅身份資訊，冒充你去營業廳補辦一張新嘅 SIM 卡。一旦補辦成功，你手機嘅訊號會突然消失，而攻擊者嘅手機插住嗰張新卡，可以收到你所有嘅短訊驗證碼——支付寶、銀行、微信，全部門戶大開。

\begin{tcbinfo}{點樣發現？SIM swap 嘅三個預警信號}
\begin{itemize}
  \item 手機突然顯示「無服務」/「無 SIM 卡」，但你明明冇鬱過張卡，重啟都冇用——呢個往往係被劫持嘅瞬間。
  \item 收到電信商發嚟嘅「補卡成功」或者「業務辦理成功」短訊——但你本人冇辦過。
  \item 突然收唔到驗證碼，但其他短訊都收唔到——因為短訊被劫持咗去攻擊者嘅卡上。
\end{itemize}

\end{tcbinfo}

\begin{tcbwarn}{防護措施——將 SIM 卡都「上鎖」}
\begin{itemize}
  \item 設定 SIM 卡 PIN 碼：路徑 設定 > 電話 > SIM 卡 PIN 碼。開啟之後，每次手機重啟或者換卡，都要輸 PIN 先用得到呢張卡。攻擊者就算拎到卡都用唔到。
  \item 開啟電信商二次驗證：登入中國移動/聯通/電信 App，搵到「安全中心」或者「業務辦理驗證」，開啟非本人補卡短訊提醒同業務辦理二次密碼驗證。咁樣人哋去補你嘅卡時，營業員會要求額外嘅身份核驗。
  \item 重要帳號唔好只用短訊驗證：銀行、Apple ID 等高價值帳號，優先用認證器 App（Google Authenticator、微軟 Authenticator）或者硬件密鑰，而唔係短訊驗證碼。
  \item 身份證複印件加水印：辦理業務需要嘅複印件寫明「僅供 XX 業務使用」——攻擊者往往靠洩露嘅身份證資訊冒充你。
\end{itemize}

\end{tcbwarn}

\textbf{被劫持後嘅應急處理——速度就係一切}\par

\begin{enumerate}
  \item 即刻撥打電信商客服掛失：移動 10086 / 聯通 10010 / 電信 10000，要求即刻停機並凍結補卡業務，阻止攻擊者繼續使用你嘅號碼。
  \item 用其他裝置登入關鍵帳號：用電腦或者家人手機，即刻修改銀行、支付寶、微信、Apple ID 嘅密碼，並踢走陌生裝置。
  \item 凍結支付帳戶：支付寶 95188 掛失，微信安全中心凍結，銀行客服凍結網銀。
  \item 帶身份證去營業廳：本人持身份證去自營廳補辦新卡，恢復號碼使用權。
  \item 報警：SIM swap 已構成「侵犯公民個人資訊罪」，帶齊證據（停機時間記錄、資金損失、補卡記錄）報案。
  \item 全面排查：恢復號碼後，逐一檢查近期登入記錄同資金流水，確認有冇被冒用過嘅痕跡。
\end{enumerate}

\begin{tcbwarn}{預防勝於治療——日常 5 件事}
\begin{enumerate}
  \item 鎖屏密碼 + 雙重認證，兩個都開
  \item 「查找我嘅手機」確認開啟
  \item SIM 卡設定 PIN 碼（設定 > 電話 > SIM 卡 PIN 碼），唔見咗手機人哋唔可以用你嘅卡
  \item 唔安裝來歷不明嘅 App，唔點陌生連結
  \item 定期備份重要數據（相片、通訊錄、聊天記錄）
\end{enumerate}

\end{tcbwarn}

\begin{tcbpractice}{動手：做一份屬於你嘅「應急聯絡卡」}
\begin{enumerate}
  \item 搵一張紙，寫低以下號碼：電信商客服（移動 10086/聯通 10010/電信 10000）、銀行客服、96110 反詐專線、110 報警電話、家人電話。
  \item 寫低你嘅 Apple ID/品牌帳號嘅註冊電郵同備用聯絡方式。
  \item 將呢張紙放喺銀包入面或者貼喺雪櫃上面——唔好存喺手機入面，因為出事時手機可能唔喺身邊。
  \item 而家就檢查：「查找我嘅手機」係咪已開啟，SIM 卡 PIN 碼係咪已設定。
\end{enumerate}

\end{tcbpractice}

\begin{tcbsummary}{本章小結}
\begin{itemize}
  \item 手機唔見咗黃金 30 分鐘：藉手機撥打遙距鎖定掛失 SIM 卡凍結支付通知親友報警
  \item 俾人呃咗之後 72 小時：停止轉帳保留證據撥 96110報警聯絡銀行改密碼
  \item 手機中毒：安全模式卸載可疑 App安全掃描復原出廠（終極方案）
  \item 帳號被盜：改密碼踢裝置查登入記錄開雙重認證查綁定查支付
  \item 勒索訊息 99.9\% 係詐騙，唔回覆唔畀錢，截圖拉黑
  \item AI 換臉/擬聲詐騙：睇眨眼/嘴型/邊緣/光影破綻，叫對方轉頭掩面、問私密問題、回撥已知號碼
  \item SIM swap 劫持：手機突然「無服務」即警惕；設 SIM 卡 PIN、開電信商二次驗證、重要帳號改用認證器 App
  \item 日常預防五件事：鎖屏+2FA、查找手機、SIM 卡 PIN、唔裝來歷不明 App、定期備份
\end{itemize}

\end{tcbsummary}
